% !TEX root = ../Thesis.tex

\section{What Is Estimated, and Why It Is Not Validated Like a Prediction}

The object of interest is the individual treatment effect
$\tau(x) = \mathbb{E}[Y(1) - Y(0) \mid X = x]$, the difference between the two
potential outcomes of the same patient, one under each treatment. Of this
difference only one term is ever observed, because every patient receives
exactly one of the two treatments: the other outcome, the one that would have
occurred under the treatment not given, does not exist in the data for anyone.
This isthe fundamental problem of causal inference, and it is not a data-quality
issue that a larger cohort would fix.

A direct consequence of this estimantions is their uncomparability, there is nothing
to compare them against. A possible way of verifying the estimator is on the aggregate,
by comparing the treatment effect across subgroups of patients 
for which the missing counterfactual reappears in the form of a group average estimable 
from the patients who did receive the treatment in question. Every validation strategy
used in this thesis is a consequence of this one constraint.

\section{Five Families of Estimators, Chosen to Err Differently}
\label{sec:five-estimators}

No single method is used, because a result that depends on the internal
assumptions of one algorithm is not distinguishable from an artefact of that
algorithm. Five estimator families are used instead, deliberately spanning
different ways of failing, so that a conclusion which survives all five stops
being attributable to any one of them.

\paragraph{Two-model meta-learner (T-learner).}
The simplest of the five: fit one outcome model per treatment arm,
$\hat\mu_1(x)$ and $\hat\mu_0(x)$, and take their difference as
$\hat\tau(x) = \hat\mu_1(x) - \hat\mu_0(x)$. Its strength is that any
regression or classification method can be reused as a component; its weakness
is the same fact seen from the other side — it inherits every bias of the two
component models, and small independent errors in each one do not cancel in
the subtraction, they add.

Two consequences follow for how $\hat\mu_1$ and $\hat\mu_0$ are built. Both
arms need the same model family, tuned separately on each arm's data:
differencing two algorithms with different inductive biases, a random forest
against a logistic regression for instance, produces a surface whose variation
reflects those biases as much as any treatment effect, and at the level of the
output that spurious variation is indistinguishable from heterogeneity. Both
arms' outcome probabilities also need to be calibrated. Calibration is a
monotone transformation, so it leaves each model's discrimination and its
ranking of patients unchanged, leaving classification performance unaffected.
A difference is affected: a distortion common to both arms cancels out of a
ranking metric such as AUROC but not out of a subtraction, so two
uncalibrated, well-discriminating models can still produce an apparent effect
where the true probabilities agree.

\paragraph{Causal forest.}
The causal forest estimates the treatment effect $\tau(x)$ directly,
rather than fitting two separate outcome models and taking their
difference. It uses double machine learning to account for confounding:
models for $\mathbb{E}[Y \mid X]$ and $\mathbb{E}[T \mid X]$ are first
fitted and their effects are removed before estimating the treatment
effect.

This helps prevent variables that influence treatment assignment from
being incorrectly interpreted as effect modifiers. The forest also uses
\emph{honest} splitting: one subset of the data is used to determine the
tree structure, while a separate subset is used to estimate the
treatment effects within each leaf. This reduces the risk of finding
patterns that are specific to the data used to build the tree.

\paragraph{Bayesian causal forest (BCF).}
A Bayesian counterpart that separates the prognostic surface, how sick a
patient already is regardless of treatment, from the effect surface,
$\tau(x)$ itself, and puts a separate prior and a separate tree ensemble on
each. Because a stronger, more differentiated prior can be placed on the
prognostic side, a residual confounder that produces a large main effect is
less able to masquerade as heterogeneity; the price is that the effect
surface can, in turn, be dominated by its own regularization if that prior is
too tight. It returns a full posterior over $\tau(x)$ rather than a point
estimate, which is what later sections use to ask not just what the effect
is but how sure the model is about it.

\paragraph{In-context model.}
Rather than being fitted to one dataset, this estimator is pre-trained once on
a large distribution of synthetic causal problems and, at inference time,
conditions on the new cohort without any further fitting — the same mechanism
that makes a language model adapt to a prompt it has never seen. Its errors
therefore have a different origin from the other four: not sampling variance
from a specific, finite cohort, but a domain bias inherited from the
distribution of synthetic problems it was pre-trained on. That difference in
kind is what makes it a genuinely independent check rather than one more
variant of the same idea.

\vspace{10pt}

\paragraph{Interaction forest (S-learner).}
Fits a single classification forest on the augmented design
$[X,,T,,X!\cdot!T]$ — covariates, treatment, and their interactions — and
recovers the implied CATE by toggling the treatment column and differencing
the two predictions. Because both treatment arms are fitted jointly,
information is pooled across them, which can help stabilize estimation in
the arm with fewer events. This pooling, however, can limit the detection of
treatment-effect heterogeneity: interaction effects must compete for splits
with the treatment's main effect and the prognostic signal in $X$, so weak
interactions may be attenuated toward zero rather than clearly identified.

\vspace{10pt}

None of the five is included because it is the best estimator in the
abstract — under different data each could plausibly outperform the others.
They are included because their failure modes do not overlap: two-model
error compounding, a forest's limited splitting language, a Bayesian model's
own regularization, a pre-trained model's domain bias, and an interaction
model's shrinkage toward zero are five different ways of being wrong, and an
agreement that survives all five is no longer explainable by any single one
of them.

\section{The Ranking as the Object of Validation}

Aggregation gives back what the missing counterfactual takes away, but only
once there are groups to aggregate over, and nothing in the data defines
them. The estimator itself supplies the grouping: ordering patients by their
predicted $\tau(x)$ turns a quantity attached to individuals into a family of
nested subgroups — the top 5\%, the top 10\%, and so on — each of which has a
group average treatment effect that is estimable. What is tested is then not
the value assigned to any patient but a property of the ordering: the
patients the model places at the top must, as a group, actually show a larger
effect than those it places below. How large that effect is, is a separate
question, asked and answered separately.

\vspace{10pt}

The group averages are built from a \emph{doubly robust} pseudo-outcome, one
number per patient, formed by taking the prediction of an outcome model and
correcting it with the inverse-propensity-weighted residual of the arm the
patient was actually assigned to. Individually this number is far too noisy
to be read as that patient's treatment effect — it is not an estimate of
$\tau(x)$ and is never used as one — but its average over any subgroup is a
consistent estimate of that subgroup's average effect, and it stays
consistent if either the outcome model or the treatment-assignment model is
correctly specified. Here the second is correct by construction: treatment
was allocated at random, so the assignment probability is known rather than
estimated, and the validity of every group average that follows does not rest
on the outcome model being right.

\begin{equation}
\hat{\varphi}_i \;=\; \hat{\mu}(1, X_i) - \hat{\mu}(0, X_i)
\;+\; \frac{T_i - \hat{e}(X_i)}{\hat{e}(X_i)\big(1 - \hat{e}(X_i)\big)}
\Big(Y_i - \hat{\mu}(T_i, X_i)\Big)
\label{eq:aipw-pseudo-outcome}
\end{equation}

\noindent where
\begin{description}
    \item[$X_i$] covariate vector for patient $i$;
    \item[$T_i \in \{0,1\}$] treatment actually received by patient $i$ (e.g. $1$ = prolonged DAPT, $0$ = abbreviated DAPT);
    \item[$Y_i$] observed outcome for patient $i$;
    \item[$\hat{\mu}(t, x)$] outcome model prediction for treatment arm $t$ given covariates $x$;
    \item[$\hat{e}(x)$] propensity score, $P(T=1 \mid X=x)$; known by design under randomization rather than estimated;
    \item[$\hat{\varphi}_i$] doubly robust AIPW pseudo-outcome for patient $i$, whose subgroup average is a consistent estimator of the subgroup's average treatment effect.
\end{description}

\vspace{10pt}

The \textbf{TOC curve} (targeting operator characteristic) reads the ordering
off at every depth. For each fraction $q$ it takes the patients ranked
highest and reports the difference between the average effect within that
group and the average effect over the whole cohort — equivalently, the gain
over selecting a group of the same size at random. A ranking that carries no
information gives a curve flat at zero, and every curve is pinned to zero at
$q=1$, where the selected group is the cohort itself; a ranking that works
gives a curve that is positive and decreasing, since concentrating on fewer
patients should concentrate the effect. The thresholds that cut the ranking
are fixed on the data used to fit the estimator and then applied to held-out
patients, whose pseudo-outcomes supply the effects, so no part of the curve
is read off the sample that produced the ranking being tested.

\vspace{10pt}

A second test asks the question the ordering deliberately sets aside.
Grouping patients into quantiles of predicted CATE, \emph{group calibration}
compares the effect actually estimated within each group against the effect
the model predicts for that same group, and scales the discrepancy by how far
the group effects lie from the overall average effect. The resulting
calibration $R^2$ falls to zero or below exactly when the model does no
better than assigning every patient the cohort average. The two tests
therefore fail in different ways: a model can be well calibrated and still
order nobody usefully, if the effect it correctly predicts is the same for
everyone, and it can order patients correctly while overstating by a constant
factor how much they differ.

\vspace{10pt}

Two scalar summaries integrate the TOC curve, differing only in how they
weight its depths. \textbf{AUTOC} weights every fraction equally, so a narrow
high-benefit group at the top contributes as much as a broad one; it is the
more powerful summary when the benefit is concentrated in a small subgroup.
\textbf{QINI} weights each point by the size of the group it refers to, which
suppresses the top of the ranking and rewards an effect spread broadly across
the population. Which of the two is the appropriate lens depends on whether
the responsive subgroup is a narrow tail or a large fraction of the cohort —
precisely what is not known before the analysis — so both are reported, and
only a conclusion that holds under both weightings is a property of the
ranking rather than of the choice of weight.

\section{The Noise Floor}

A positive AUTOC, or any nonzero dispersion of the estimated $\tau(x)$ across
patients, is not by itself evidence of heterogeneity: an estimator applied to
a constant true effect still produces a spread of estimates, because
estimation error alone has variance, and that spread grows with how flexible
the model is allowed to be, independently of whether the underlying effect
varies at all. Comparing against zero therefore answers the wrong question.

\vspace{10pt}

The comparison used throughout this thesis is against a \textbf{noise floor}
obtained by permuting the treatment labels and repeating the identical
estimation and evaluation pipeline. Permutation destroys any real treatment
effect while leaving the covariates, the outcome rates, and every modelling
choice untouched, so the resulting distribution of AUTOC — or of CATE
dispersion — is exactly what the pipeline manufactures from nothing. A result
is treated as a signal only if it clears this floor, not if it merely clears
zero.

\section{From Estimation to Decision: The Value of a Rule}

A ranking can be statistically correct and still not be worth acting on if the
gain it identifies is too small to justify changing the treatment strategy.
The next step therefore asks a different question: not whether the model ranks
patients correctly, but whether \emph{following} that ranking leads to a
meaningful improvement in outcomes.

\vspace{10pt}

This requires a treatment policy that determines which treatment should be
given to each patient based on the expected balance of gains and losses.
Ranking patients using only one selected outcome may not be sufficient, since
improving one outcome can come at the cost of worsening another. The policy
must therefore consider the relevant outcomes jointly and evaluate whether
the resulting treatment decisions provide an overall benefit.

\vspace{10pt}

For each candidate policy, its expected value is estimated using a
cross-fitted AIPW (augmented inverse-propensity-weighted) estimator. The
estimated value is then compared with simple fixed strategies that do not use
the model, such as treating everyone with the same treatment. Confidence
intervals are used to assess whether the difference is meaningful. This
changes the question from \emph{`Does the model rank patients correctly?''}
to \emph{`Does following its recommendations actually improve outcomes?''}


\section{The Consequence for Model Selection}

If the quantity that matters is the quality of the ranking, the
hyperparameters of every estimator must be chosen against that same target,
not against a prediction-error criterion that has no direct bearing on it.
Tuning therefore optimizes the cross-validated \emph{lower bound} of the
AUTOC, $\text{AUTOC} - z \cdot \text{SE}$ averaged across folds, rather than
the point estimate. The choice is a deliberate defence against the winner's
curse: across many configurations tried, the one that looks best often looks
so partly by chance, and a criterion that rewards a high score obtained with
a large standard error would systematically select the luckiest configuration
rather than the best one.

\vspace{10pt}

A second defence handles a specific failure mode directly: a configuration
that collapses to a constant CATE for every patient is not discarded as
invalid, it is scored zero. Without this rule such a configuration could hide
behind whatever validation score it happens to draw by chance, since a
model with no ranking at all cannot be wrong about one either. The same
tuning procedure searches over the flexibility of the nuisance models jointly
with the flexibility of the effect model, because a less noisy doubly-robust
pseudo-outcome raises the achievable lower bound exactly as much as a more
expressive forest does, and optimizing one while holding the other fixed
would credit the wrong component for the gain.

\vspace{10pt}

The result is that the measure used to validate an estimator, the criterion
used to select among tuning configurations, and the criterion used to accept
or reject a claim of heterogeneity are, by construction, the same quantity.
No separate notion of accuracy is allowed to override what the ranking-based
evaluation of Section~2.3 says.
